\documentclass[serif,mathserif,final]{beamer}
\mode<presentation>{\usetheme{DFTFE}}
\usepackage{amsmath,amsfonts,amssymb,pxfonts,eulervm,xspace}
\usepackage{graphicx}
\usepackage{ragged2e}
\usepackage{psfrag}
\usepackage{caption}
\usepackage{subcaption}
\usepackage{float}
\usepackage{geometry}
\usepackage[many]{tcolorbox}
\usepackage{upgreek}
%\usepackage{auto-pst-pdf}
\graphicspath{{./figures/}}
\usepackage[orientation=portrait,size=custom,height=119,width=84, scale=1.3]{beamerposter}
\usepackage{multirow}
\usepackage{braket}
\DeclareMathAlphabet{\mathpzc}{OT1}{pzc}{m}{it}
\newcommand{\DFTFE}{\texttt{DFT-FE}}
\newcommand{\ASE}{\texttt{ASE}}

%-------------------------------------------------------------------
\title{\LARGE Daemonized DFT-FE\\
\vspace{0.3em}
\large A persistent, integration-ready DFT platform for scalable innovation in science and industry}
\author{Mehul Darak, Kartick Ramakrishnan, Prof. Phani Motamarri (PI)}
\institute[]{\normalsize{\color{black}Department of Computational and Data Sciences, Indian Institute of Science, Bengaluru, India\color{black}}}
%-------------------------------------------------------------------

% (keep all your macro definitions exactly as you had them)
\newcommand{\parder}[2]{\frac{\partial #1}{\partial #2}}
\newcommand{\parsecder}[3]{\frac{\partial^2 #1}{\partial #2\; \partial #3}}
\newcommand{\shp}[2]{N^{#1}_{#2}(\bx)}
\newcommand{\bzero}{\mathbf{0}}
\newcommand{\del}{\nabla}
\newcommand{\bPsi}{\boldsymbol{\Psi}}
\newcommand{\bvPsi}{\boldsymbol{\varPsi}}
\newcommand{\varphibar}{\bar{\varphi}}
\newcommand{\bcolon}{\boldsymbol{:}}
\newcommand{\bb}{\boldsymbol{b}}
\newcommand{\bc}{\boldsymbol{c}}
\newcommand{\bff}{\boldsymbol{f}}
\newcommand{\bg}{\boldsymbol{g}}
\newcommand{\bh}{\boldsymbol{h}}
\newcommand{\bk}{\boldsymbol{k}}
\newcommand{\bl}{\boldsymbol{l}}
\newcommand{\bn}{\boldsymbol{n}}
\newcommand{\bp}{\boldsymbol{p}}
\newcommand{\bq}{\boldsymbol{q}}
\newcommand{\br}{\boldsymbol{r}}
\newcommand{\bv}{\boldsymbol{v}}
\newcommand{\bs}{\boldsymbol{s}}
\newcommand{\bt}{\boldsymbol{t}}
\newcommand{\bu}{\boldsymbol{u}}
\newcommand{\bx}{\boldsymbol{\textbf{x}}}
\newcommand{\bA}{\boldsymbol{A}}
\newcommand{\bB}{\boldsymbol{B}}
\newcommand{\bC}{\boldsymbol{C}}
\newcommand{\bE}{\boldsymbol{E}}
\newcommand{\bF}{\boldsymbol{F}}
\newcommand{\bH}{\boldsymbol{\textbf{H}}}
\newcommand{\bI}{\boldsymbol{I}}
\newcommand{\bK}{\boldsymbol{K}}
\newcommand{\bKbar}{\bar{\bK}}
\newcommand{\bxbar}{\bar{\boldsymbol{x}}}
\newcommand{\bffbar}{\bar{\bff}}
\newcommand{\bX}{\boldsymbol{X}}
\newcommand{\bM}{\boldsymbol{\textbf{M}}}
\newcommand{\bN}{\boldsymbol{N}}
\newcommand{\bO}{\boldsymbol{O}}
\newcommand{\bP}{\boldsymbol{P}}
\newcommand{\bR}{\boldsymbol{R}}
\newcommand{\bS}{\boldsymbol{S}}
\newcommand{\bT}{\boldsymbol{T}}
\newcommand{\bU}{\boldsymbol{U}}
\newcommand{\bW}{\boldsymbol{W}}
\newcommand{\rhoTilde}{\widetilde{\rho}}
\newcommand{\bphi}{\boldsymbol{\phi}}
\newcommand{\bLam}{\boldsymbol{\Lambda}}
\newcommand{\tbar}{\bar{\bt}}
\newcommand{\ubar}{\bar{\bu}}
\newcommand{\ubarn}{\bar{u}}
\newcommand{\uhat}{\hat{\bu}}
\newcommand{\fhat}{\hat{\bff}}
\newcommand{\ghat}{\hat{\bg}}
\newcommand{\phihat}{\hat{\bphi}}
\newcommand{\lambar}{\bar{\lambda}}
\newcommand{\phibar}{\bar{\phi}}
\newcommand{\psibar}{\bar{\psi}}
\newcommand{\Psibar}{\bar{\Psi}}
\newcommand{\rhobar}{\bar{\rho}}
\newcommand{\rhoin}[2]{\rho_{\text{in}}^{{#1}^{(#2)}}}
\newcommand{\rhoout}[2]{\rho_{\text{out}}^{{#1}^{(#2)}}}
\newcommand{\rhoinout}[2]{\rho_{\text{in(out)}}^{{#1}^{(#2)}}}
\newcommand{\epsilonbar}{\bar{\epsilon_{i}}}
\newcommand{\mubar}{\bar{\mu}}
\newcommand{\bvPhi}{\boldsymbol{\varPhi}}
\newcommand{\Vbar}{\bar{V}}
\newcommand{\psin}{\widetilde{\psi}_n}
\newcommand{\Psin}{\widetilde{\bPsi}}
\newcommand{\wbar}{\bar{w}}
\newcommand{\intomega}{\int_{\varOmega}}
\newcommand{\domega}{\;d\varOmega}
\newcommand{\dr}{\,d\br}
\newcommand{\dx}{\,d\bx}
\newcommand{\btH}{\boldsymbol{\widetilde{\textbf{H}}}}
\newcommand{\tN}{\widetilde{N}}
\definecolor{hellgruen}{rgb}{0.2,0.7,0.2}
\newcommand{\cb}{\color{blue}}
\newcommand{\cred}{\color{red}}
\newcommand{\cn}{\color{black}}
\newcommand{\cm}{\color{magenta}}
\newcommand{\cbh}{\color{hellgruen}}
\newcommand{\bPhi}{\boldsymbol{\Phi}}
\newcommand{\chieps}{\,\chi_{\varepsilon}(\bx)\,}
\newcommand{\bxprime}{\boldsymbol{\textbf{x}^{\prime}}}
\newcommand{\bxdprime}{\boldsymbol{\textbf{x}^{\prime\prime}}}
\newcommand{\bdir}{\boldsymbol{\Upsilon}}
\newcommand{\atzero}{\bigg|_{\varepsilon = 0}}
\newcommand{\derveps}{\frac{d}{d\varepsilon}}
\newcommand{\innerproductcomma}[2]{\langle #1, #2 \rangle}
\usepackage{makecell}
\newcommand{\order}{\mathcal{O}}
\renewcommand\theadalign{bc}
\renewcommand\theadfont{\bfseries}
\renewcommand\theadgape{\Gape[4pt]}
\renewcommand\cellgape{\Gape[4pt]}

\usepackage{tikz}
\usetikzlibrary{arrows.meta,decorations.pathreplacing,arrows,decorations.markings,positioning,shapes,matrix}
\tikzstyle{vecArrow} = [thick, decoration={markings,mark=at position
   1 with {\arrow[semithick]{open triangle 60}}}, double distance=1.4pt, shorten >= 5.5pt,
   preaction = {decorate}, postaction = {draw,line width=1.4pt, white,shorten >= 4.5pt}]
\tikzset{mynode/.style={rectangle,rounded corners,draw=black, top color=white, bottom color=yellow!50,very thick, inner sep=0.5em, minimum size=1em, text centered},
         myarrow/.style={->, >=latex', shorten >=1pt, thick},
         mylabel/.style={text width=7em, text centered}}
\tikzset{myarrow/.style={draw, fill=orange, single arrow, minimum height=4.5ex, single arrow head extend=0.5ex}}
\tikzset{myarrowsmall/.style={draw, fill=orange, single arrow, minimum height=4.5ex, single arrow head extend=0.25ex}}
\newcommand{\arrowup}{\tikz [baseline=-0.5ex]{\node [myarrow,rotate=90] {};}}
\newcommand{\arrowdown}{\tikz [baseline=-1ex]{\node [myarrow,rotate=-90] {};}}
\newcommand{\arrowright}{\tikz [baseline=-1ex]{\node [myarrow] {};}}
\newcommand{\arrowrightsmall}{\tikz [baseline=-1ex]{\node [myarrowsmall] {};}}
\newtcolorbox{myboxii}[1][]{breakable,freelance,title=#1,colback=white,
  colbacktitle=white,coltitle=black,fonttitle=\bfseries,bottomrule=0pt,boxrule=0pt,colframe=white,
  overlay unbroken and first={\draw[red!75!black,line width=3pt]([xshift=5pt]frame.north west)--(frame.north west)--(frame.south west);
  \draw[red!75!black,line width=3pt]([xshift=-5pt]frame.north east)--(frame.north east)--(frame.south east);},
  overlay unbroken app={\draw[red!75!black,line width=3pt,line cap=rect](frame.south west)--([xshift=5pt]frame.south west);
  \draw[red!75!black,line width=3pt,line cap=rect](frame.south east)--([xshift=-5pt]frame.south east);},
  overlay middle and last={\draw[red!75!black,line width=3pt](frame.north west)--(frame.south west);
  \draw[red!75!black,line width=3pt](frame.north east)--(frame.south east);},
  overlay last app={\draw[red!75!black,line width=3pt,line cap=rect](frame.south west)--([xshift=5pt]frame.south west);
  \draw[red!75!black,line width=3pt,line cap=rect](frame.south east)--([xshift=-5pt]frame.south east);} }

%========================
\begin{document}
% Bigger block titles and bodies
\setbeamerfont{block title}{size=\Large,series=\bfseries}
\setbeamerfont{block body}{size=\large}
% Optional: make normal text a tad bigger too
\setbeamerfont{normal text}{size=\large}

%========================
% PAGE 1
%========================
%========================
% PAGE 1 — FULL POSTER
%========================
\begin{frame}[t]
    %--- Header ---
    \begin{beamercolorbox}[sep=1em,center,wd=\paperwidth]{postertitle}
        \usebeamerfont{title}{\huge \textbf{Daemonized DFT-FE: A Persistent Interface for Scalable ab initio Workflows}}\\[0.5em]
        \usebeamerfont{author}{\Large Mehul Darak, Kartick Ramakrishnan, Phani Motamarri (phanim@iisc.ac.in)}\\[0.3em]
        \usebeamerfont{institute}{\large Department of Computational and Data Sciences, Indian Institute of Science, Bengaluru, India}
    \end{beamercolorbox}
    \vspace{1cm}

    %--- Columns ---
    \begin{columns}[t]

        %========================
        % LEFT COLUMN
        %========================
        \begin{column}{0.48\paperwidth}

            %--- BLOCK 1: Motivation ---
            \begin{tcolorbox}[title=Motivation: The I/O Bottleneck]
                \justifying
                Traditional couplings between density functional theory (DFT) solvers and workflow engines (e.g., \ASE{}) rely on file-based communication. [cite_start]This approach introduces significant inefficiencies for high-throughput iterative workflows[cite: 7, 18]:
                \begin{itemize}
                    \item \textbf{Process Overhead:} Relaunches a new DFT process for every energy/force evaluation.
                    \item \textbf{I/O Latency:} Repeatedly reloads pseudopotentials and writes intermediate files.
                    \item \textbf{Data Waste:} Discards valuable wavefunction data, forcing "cold starts" at every step.
                \end{itemize}
                
                \textbf{The Solution:} We introduce a \textbf{socket-driven persistent interface} for \DFTFE{}, a massively parallel real-space finite-element DFT code. [cite_start]This transforms \DFTFE{} into a continuously running service[cite: 11].
            \end{tcolorbox}
            \vspace{1em}

            %--- BLOCK 2: Architecture ---
            \begin{tcolorbox}[title=System Architecture \& Data Flow]
                \justifying
                [cite_start]The architecture follows a Client-Server (Driver-Daemon) model adhering to the \textbf{i-PI communication protocol}[cite: 14, 45].
                
                \begin{figure}
                    \centering
                    % PLACEHOLDER FOR FIGURE 1 from Manuscript
                    %                     % Ensure you have the image file in the directory
                    \includegraphics[width=0.95\linewidth]{figure1.png} 
                    \caption{\textbf{Persistent Service Mode:} The daemon initializes once. [cite_start]For each step, it receives atomic structures, performs the SCF solve reusing previous states, and returns $\{E, \mathbf{F}, \boldsymbol{\sigma}\}$[cite: 30, 31].}
                \end{figure}

                \begin{itemize}
                    [cite_start]\item \textbf{Initialization:} MPI setup, mesh generation, and pseudopotential loading occur only once[cite: 36].
                    \item \textbf{Loop:} Listen $\rightarrow$ Compute $\rightarrow$ Respond.
                    [cite_start]\item \textbf{Protocol:} JSON messages over lightweight sockets[cite: 46].
                \end{itemize}
            \end{tcolorbox}
            \vspace{1em}

            %--- BLOCK 3: ASE Client Integration ---
            \begin{tcolorbox}[title=ASE Integration: \texttt{DFTFESocketCalculator}]
                \justifying
                [cite_start]The Python client seamlessly replaces file-based workflows with zero code modification for the end-user[cite: 52, 63].
                
                \textbf{Example Implementation:}
                \begin{lstlisting}[language=Python]
from ase.io import read
from ase.optimize import BFGS
from dftfe_socket import DFTFESocketCalculator

# Define Calculator
calc = DFTFESocketCalculator(
    command="srun -n 64 dftfe-socket", # HPC Launch
    xc="PBE", kpts=(2,2,2), h=0.25     # Static Params
)

atoms = read("structure.cif")
atoms.calc = calc

# Geometry Optimization (Persistant Session)
with calc:
    dyn = BFGS(atoms)
    dyn.run(fmax=0.02) # Iterative calls over socket
                \end{lstlisting}
                [cite_start]Parameters like XC functional and basis size are set once; coordinates are transmitted dynamically[cite: 58].
            \end{tcolorbox}

        \end{column}

        %========================
        % RIGHT COLUMN
        %========================
        \begin{column}{0.48\paperwidth}

            %--- BLOCK 4: HPC Methodology & Parallelism (HiPC Theme) ---
            \begin{tcolorbox}[title=HPC Methodology \& Scalability]
                \justifying
                \DFTFE{} is designed for scalability on large HPC systems using Finite-Element discretization. [cite_start]The socket interface preserves this massive parallelism[cite: 9, 64].
                
                \textbf{Algorithmic Implementation:}
                \begin{itemize}
                    \item \textbf{Parallel Execution:} Users launch via standard schedulers (e.g., \texttt{srun -n 256}). [cite_start]The socket mode maintains full MPI parallelism[cite: 66, 70].
                    [cite_start]\item \textbf{Domain Decomposition:} While Rank 0 manages external socket I/O, all ranks participate in the SCF calculation through \DFTFE{}'s native domain decomposition[cite: 67].
                    [cite_start]\item \textbf{Communication Efficiency:} The socket overhead is negligible compared to the SCF computation costs for large systems[cite: 68].
                    [cite_start]\item \textbf{Fault Tolerance:} The persistent connection allows for graceful error handling (e.g., SCF non-convergence) without crashing the pipeline, enabling dynamic parameter adjustment[cite: 73].
                \end{itemize}
            \end{tcolorbox}
            \vspace{1em}

            %--- BLOCK 5: Computational Efficiency ---
            \begin{tcolorbox}[title=Accelerating Convergence: Warm Starts]
                \justifying
                A critical advantage of the daemon mode for Materials Science is the ability to maintain state in memory.
                
                \begin{itemize}
                    [cite_start]\item \textbf{Wavefunction Reuse:} In iterative workflows like MD or relaxation, atomic displacements between steps are small[cite: 42].
                    [cite_start]\item \textbf{SCF Speedup:} By reusing the density and wavefunctions from the previous step as the initial guess, the number of SCF iterations required for convergence is substantially reduced compared to cold restarts[cite: 13, 43].
                    [cite_start]\item \textbf{I/O Reduction:} Eliminates filesystem overheads in multi-node environments where repeated read/write operations dominate runtime[cite: 72].
                \end{itemize}
            \end{tcolorbox}
            \vspace{1em}
            
            %--- BLOCK 6: Applications & Conclusion ---
            \begin{tcolorbox}[title=Impact \& Applications, colframe=accentred]
                \justifying
                [cite_start]This full-stack integration positions \DFTFE{} as a high-throughput backend for modern scientific workflows[cite: 75, 80].
                
                \textbf{Key Applications:}
                \begin{itemize}
                    \item \textbf{Ab Initio Molecular Dynamics (AIMD):} Long-timescale simulations.
                    [cite_start]\item \textbf{Machine Learning:} On-the-fly training of interatomic potentials and active learning[cite: 78].
                    \item \textbf{Structure Search:} Large-scale geometry optimizations.
                \end{itemize}
                
                [cite_start]\textbf{Conclusion:} The socket interface eliminates redundant initialization and disk I/O while preserving \DFTFE{}'s strong scaling across thousands of cores, bridging the gap between traditional HPC solvers and Python-based automation[cite: 76, 78].
            \end{tcolorbox}
            \vspace{1em}

            %--- References & Ack ---
            \begin{beamercolorbox}[wd=\linewidth, leftskip=1em]{references}
                \small \textbf{References:}
                [1] ASE Project, "Communication with calculators over sockets."
                [2] Ceriotti et al., Comput. Phys. Commun., 2014.
                [3] Motamarri et al., "DFT-FE 1.0", Comput. Phys. Commun., 2022.
                
                \vspace{0.5em}
                [cite_start]\textbf{Acknowledgments:} The authors acknowledge the support of IISc's MATRIX Lab[cite: 82].
            \end{beamercolorbox}

        \end{column}
    \end{columns}
\end{frame}
\end{document}
